\documentclass[a4paper]{article}

% Shared Preamble for MathLatexDocs

% --- Core Packages ---
\usepackage[english]{babel}
\usepackage[utf8]{inputenc}
\usepackage{amsmath}
\usepackage{amsthm}
\usepackage{amssymb}
\usepackage{mathtools}
\usepackage{graphicx}
\usepackage{microtype} % Better justification and fewer overfull hboxes
\usepackage{bm}
\usepackage{enumitem}
\usepackage{booktabs}
\usepackage{mathrsfs}

% --- Layout & Design ---
\usepackage{docmute}
\usepackage{subfiles} % Support for master/subfile structure
\usepackage[colorinlistoftodos]{todonotes}
\usepackage[colorlinks=true, linkcolor=blue, citecolor=magenta]{hyperref}
\usepackage{cleveref} % Better cross-referencing
\usepackage{tcolorbox} % For styled boxes
\usepackage{tikz} % For drawing diagrams
\usetikzlibrary{arrows.meta,positioning}

% --- Theorem Environments (Classic Style) ---
\makeatletter
\@ifundefined{classicthm}{%
  \newtheorem{classicthm}{Theorem}[section]
  \newtheorem{classicdef}[classicthm]{Definition}
  \newtheorem{classiclem}[classicthm]{Lemma}
  \newtheorem{theorem}{Theorem}[section]
  \newtheorem{corollary}{Corollary}[theorem]
  \newtheorem{lemma}[theorem]{Lemma}
  \newtheorem{definition}{Definition}[section]
  \newtheorem{proposition}[theorem]{Proposition}
  \newtheorem{remark}[theorem]{Remark}
  \newtheorem{example}[theorem]{Example}
}{}

% Wrappers to maintain compatibility with the {Title}{Label} syntax used in files
\@ifundefined{mytheorem}{%
  \newenvironment{mytheorem}[2]{%
    \begin{classicthm}[#1]\label{thm:#2}
  }{%
    \end{classicthm}
  }
  
  \newenvironment{mydefinition}[2]{%
    \begin{classicdef}[#1]\label{def:#2}
  }{%
    \end{classicdef}
  }
  
  \newenvironment{mylemma}[2]{%
    \begin{classiclem}[#1]\label{lem:#2}
  }{%
    \end{classiclem}
  }
}{}
\makeatother

% --- Custom Commands ---
\providecommand{\Z}{\mathbb{Z}}
\providecommand{\R}{\mathbb{R}}
\providecommand{\C}{\mathbb{C}}
\providecommand{\N}{\mathbb{N}}
\providecommand{\G}{\mathcal{G}}
\providecommand{\Lop}{\mathcal{L}}
\providecommand{\PP}{\mathcal{P}}
\providecommand{\dd}{\,\mathrm{d}}
\providecommand{\bitand}{\mathbin{\&}}
\providecommand{\CFDots}{\genfrac{}{}{0pt}{0}{}{\ddots}}


\title{\textbf{The Game of Nim, Bouton's Theorem, and Fibonacci Bit-Patterns}\\[6pt]
\large A Complete Mathematical Derivation of Project Euler Problem 301}
\author{Emil Kerimov}
\date{September 2026}

\begin{document}

\maketitle

\begin{abstract}
We present a rigorous, self-contained mathematical analysis and derivation of Project Euler Problem 301. By employing Bouton's classic theorem from Combinatorial Game Theory, we express the winning/losing positions of three-heap Nim via the bitwise XOR sum (Nim-sum). We analyze the specific family of configurations $(n, 2n, 3n)$ and establish an exact algebraic bridge between bitwise addition and standard arithmetic, reducing the condition $X(n, 2n, 3n) = 0$ to carryless addition $n \bitand 2n = 0$. We prove that this condition is equivalent to the absence of consecutive ones in the binary representation of $n$, and we enumerate such integers on the interval $[1, 2^{30}]$ via a combinatorial bijection with the Fibonacci sequence, arriving at the exact count $\mathcal{N}(2^{30}) = F_{32} = 2\,178\,309$.
\end{abstract}

\tableofcontents
\vspace{1em}

\section{Introduction and Problem Formulation}

\subsection{Rules of Impartial Normal-Play Nim}
Nim is a classic impartial game played by two players with a finite collection of heaps containing discrete tokens (or stones). The game operates under the \textbf{normal play convention}:
\begin{enumerate}[label=(\roman*)]
    \item Players alternate turns.
    \item On each turn, the active player chooses a single non-empty heap and removes any positive integer number of stones from it (up to the entire heap).
    \item The game ends when no stones remain in any heap.
    \item The first player unable to make a valid move (because all heaps are empty) loses; equivalently, the player who makes the final move wins.
\end{enumerate}

A game state with three heaps is denoted by the triple $(n_1, n_2, n_3) \in \mathbb{N}_0^3$, where $n_i$ represents the number of stones in the $i$-th heap.

\subsection{Combinatorial Game States: $\mathcal{P}$-Positions and $\mathcal{N}$-Positions}
In Combinatorial Game Theory (CGT), every position in a finite impartial game without cycles belongs to one of two mutually exclusive classes:
\begin{itemize}
    \item \textbf{$\mathcal{P}$-positions (Previous-player winning / Losing for the next player):} Any valid move from a state in $\mathcal{P}$ leads to a state in $\mathcal{N}$. The terminal state $(0, 0, 0)$ is in $\mathcal{P}$.
    \item \textbf{$\mathcal{N}$-positions (Next-player winning / Winning for the next player):} There exists at least one legal move from any state in $\mathcal{N}$ to some state in $\mathcal{P}$.
\end{itemize}

The characteristic evaluation function $X(n_1, n_2, n_3)$ is defined as:
\begin{equation}
X(n_1, n_2, n_3) = 
\begin{cases}
0, & \text{if } (n_1, n_2, n_3) \in \mathcal{P} \quad (\text{second player can force a win}), \\
\ne 0, & \text{if } (n_1, n_2, n_3) \in \mathcal{N} \quad (\text{first player can force a win}).
\end{cases}
\end{equation}

\subsection{Example: The Winning Strategy from Starting Position $(1, 2, 3)$}
To illustrate the mechanism of $\mathcal{P}$-positions and the winning strategy, consider the starting state $(1, 2, 3)$. We will show that $(1, 2, 3) \in \mathcal{P}$, meaning that the second player can always force a win regardless of the first player's initial move.

The first player has six possible opening moves from $(1, 2, 3)$:
\begin{enumerate}[label=\arabic*.]
    \item \textbf{Reduce Heap 1 ($1 \to 0$):} The state becomes $(0, 2, 3)$. 
    
    \textit{Winning Response by Player 2:} Player 2 reduces Heap 3 from $3 \to 2$, leaving $(0, 2, 2)$. This is a two-heap symmetric position. Whatever Player 1 takes from one heap of size 2, Player 2 mirrors by taking the exact same amount from the other heap, guaranteeing the last move.
    
    \item \textbf{Reduce Heap 2 ($2 \to 1$):} The state becomes $(1, 1, 3)$.
    
    \textit{Winning Response by Player 2:} Player 2 removes all stones from Heap 3 ($3 \to 0$), leaving $(1, 1, 0)$. Again, Player 2 adopts the mirror strategy on the two equal heaps of size 1.
    
    \item \textbf{Reduce Heap 2 ($2 \to 0$):} The state becomes $(1, 0, 3)$.
    
    \textit{Winning Response by Player 2:} Player 2 reduces Heap 3 from $3 \to 1$, leaving $(1, 0, 1)$, securing the symmetric mirror state.
    
    \item \textbf{Reduce Heap 3 ($3 \to 2$):} The state becomes $(1, 2, 2)$.
    
    \textit{Winning Response by Player 2:} Player 2 takes the single stone from Heap 1 ($1 \to 0$), leaving $(0, 2, 2)$.
    
    \item \textbf{Reduce Heap 3 ($3 \to 1$):} The state becomes $(1, 2, 1)$.
    
    \textit{Winning Response by Player 2:} Player 2 removes all stones from Heap 2 ($2 \to 0$), leaving $(1, 0, 1)$.
    
    \item \textbf{Reduce Heap 3 ($3 \to 0$):} The state becomes $(1, 2, 0)$.
    
    \textit{Winning Response by Player 2:} Player 2 reduces Heap 2 from $2 \to 1$, leaving $(1, 1, 0)$.
\end{enumerate}

In every single branch, Player 2 transitions the game to a state of two equal non-empty heaps $(k, k, 0)$, from which any move by Player 1 allows Player 2 to maintain equality or win immediately. An example sequence of optimal play is:
\begin{align*}
\text{Player 1 moves:} \quad &(1, 2, 3) \longrightarrow (1, 2, 1) \\
\text{Player 2 responds:} \quad &(1, 2, 1) \longrightarrow (1, 0, 1) \\
\text{Player 1 moves:} \quad &(1, 0, 1) \longrightarrow (0, 0, 1) \\
\text{Player 2 responds:} \quad &(0, 0, 1) \longrightarrow (0, 0, 0) \quad \text{(Player 2 wins!)}
\end{align*}
Hence, $X(1, 2, 3) = 0$.

\subsection{The Target Question}
Project Euler Problem 301 asks:
\begin{equation}
\text{For how many positive integers } n \le 2^{30} \text{ does } X(n, 2n, 3n) = 0?
\end{equation}


\section{Bouton's Theorem and the Nim-Sum}

\subsection{Definition: The Nim-Sum}
\begin{definition}[Nim-Sum]
Let $a, b \in \mathbb{N}_0$ have binary expansions:
\begin{equation}
a = \sum_{i=0}^\infty a_i 2^i, \quad b = \sum_{i=0}^\infty b_i 2^i \qquad (a_i, b_i \in \{0, 1\}).
\end{equation}
The \textbf{Nim-sum} of $a$ and $b$, denoted $a \oplus b$, is defined by bitwise addition modulo 2:
\begin{equation}
a \oplus b = \sum_{i=0}^\infty (a_i \oplus b_i) 2^i, \quad \text{where } a_i \oplus b_i = (a_i + b_i) \bmod 2.
\end{equation}
\end{definition}

The algebraic structure $(\mathbb{N}_0, \oplus)$ is an infinite elementary abelian 2-group:
\begin{enumerate}[label=(\roman*)]
    \item \textbf{Associativity and Commutativity:} $(a \oplus b) \oplus c = a \oplus (b \oplus c)$ and $a \oplus b = b \oplus a$.
    \item \textbf{Identity Element:} $a \oplus 0 = a$ for all $a \in \mathbb{N}_0$.
    \item \textbf{Self-Inversion:} $a \oplus a = 0$ for all $a \in \mathbb{N}_0$.
\end{enumerate}

\subsection{Bouton's Theorem}
\begin{theorem}[Bouton, 1901]\label{thm:nim:bouton}
In an $m$-heap game of Nim with heap sizes $(h_1, h_2, \dots, h_m) \in \mathbb{N}_0^m$, a position is a $\mathcal{P}$-position (losing for the player whose turn it is) if and only if the Nim-sum of all heap sizes is zero:
\begin{equation}
\bigoplus_{j=1}^m h_j = 0.
\end{equation}
Equivalently, $X(h_1, h_2, \dots, h_m) = \bigoplus_{j=1}^m h_j$.
\end{theorem}

\begin{proof}
Let $\mathcal{Z} = \left\{ (h_1, \dots, h_m) \in \mathbb{N}_0^m : \bigoplus_{j=1}^m h_j = 0 \right\}$. We prove by complete induction on the total stone count $S = \sum_{j=1}^m h_j$ that $\mathcal{Z}$ satisfies the axiomatic properties of $\mathcal{P}$-positions:

\begin{enumerate}
    \item \textbf{Base Case ($S = 0$):} The only state is $(0, \dots, 0)$. Its Nim-sum is $0 \oplus \dots \oplus 0 = 0 \in \mathcal{Z}$. The active player has no moves and loses, so $(0, \dots, 0) \in \mathcal{P}$.

    \item \textbf{Every move from $\mathcal{Z}$ leads outside $\mathcal{Z}$ ($\mathcal{P} \to \mathcal{N}$):}
    Let $(h_1, \dots, h_m) \in \mathcal{Z}$, so $\bigoplus_{j=1}^m h_j = 0$. A legal move strictly decreases one heap, say $h_k \to h'_k$ with $0 \le h'_k < h_k$. The new Nim-sum $V'$ is:
    \begin{equation}
    V' = \left( \bigoplus_{j \ne k} h_j \right) \oplus h'_k = \left( \bigoplus_{j=1}^m h_j \right) \oplus h_k \oplus h'_k = 0 \oplus h_k \oplus h'_k = h_k \oplus h'_k.
    \end{equation}
    Since $h'_k \ne h_k$, $h_k \oplus h'_k \ne 0$. Thus, $V' \ne 0$, meaning the new state is not in $\mathcal{Z}$.

    \item \textbf{From any state not in $\mathcal{Z}$, there exists a move into $\mathcal{Z}$ ($\mathcal{N} \to \mathcal{P}$):}
    Suppose $(h_1, \dots, h_m) \notin \mathcal{Z}$, so $V = \bigoplus_{j=1}^m h_j \ne 0$. Let $d = \lfloor \log_2 V \rfloor$ denote the most significant bit (MSB) of $V$, so bit $d$ of $V$ is 1.
    Since bit $d$ of $V$ is the parity of the $d$-th bits of all $h_j$, there exists at least one heap $h_k$ whose $d$-th bit is 1.
    
    We choose the replacement size $h'_k = h_k \oplus V$.
    \begin{itemize}
        \item For all bit indices $i > d$, the $i$-th bit of $V$ is 0, so bit $i$ of $h'_k$ matches bit $i$ of $h_k$.
        \item At bit index $d$, bit $d$ of $h_k$ is 1 and bit $d$ of $V$ is 1, so bit $d$ of $h'_k$ is $1 \oplus 1 = 0$.
    \end{itemize}
    Consequently, $h'_k < h_k$, so reducing heap $k$ to $h'_k$ is a valid Nim move. The new Nim-sum is:
    \begin{equation}
    V' = V \oplus h_k \oplus h'_k = V \oplus h_k \oplus (h_k \oplus V) = 0.
    \end{equation}
    Thus the resulting state is in $\mathcal{Z}$.
\end{enumerate}

By mathematical induction, $\mathcal{Z} = \mathcal{P}$, completing the proof.
\end{proof}


\section{Algebraic Reduction of $X(n, 2n, 3n) = 0$}

\subsection{Condition Equivalence}
For the configuration $(n, 2n, 3n)$, Theorem~\ref{thm:nim:bouton} states:
\begin{equation}
X(n, 2n, 3n) = 0 \iff n \oplus 2n \oplus 3n = 0.
\end{equation}
Using the self-inverting group property of $\oplus$:
\begin{equation}
n \oplus 2n \oplus 3n = 0 \iff n \oplus 2n = 3n.
\end{equation}
Observing that $3n = n + 2n$ in ordinary arithmetic, this becomes:
\begin{equation}\label{eq:nim:xor_sum_eq_arith_sum}
n \oplus 2n = n + 2n.
\end{equation}

\subsection{Arithmetic Addition vs.\ Bitwise XOR}

\begin{lemma}[Arithmetic and Bitwise Addition Identity]\label{lem:nim:add_xor_identity}
For any $a, b \in \mathbb{N}_0$,
\begin{equation}\label{eq:nim:add_xor_formula}
a + b = (a \oplus b) + 2(a \bitand b),
\end{equation}
where $\bitand$ denotes the bitwise AND operator.
\end{lemma}

\begin{proof}
Let $a = \sum_{i=0}^\infty a_i 2^i$ and $b = \sum_{i=0}^\infty b_i 2^i$ be the binary representations of $a$ and $b$, with $a_i, b_i \in \{0, 1\}$.
For each single bit position $i \ge 0$, we test all four possible binary assignments for $(a_i, b_i)$:

\begin{center}
\begin{tabular}{cccccc}
\toprule
$a_i$ & $b_i$ & $a_i + b_i$ & $a_i \oplus b_i$ & $a_i \bitand b_i$ & $(a_i \oplus b_i) + 2(a_i \bitand b_i)$ \\
\midrule
0 & 0 & 0 & 0 & 0 & $0 + 2(0) = 0$ \\
0 & 1 & 1 & 1 & 0 & $1 + 2(0) = 1$ \\
1 & 0 & 1 & 1 & 0 & $1 + 2(0) = 1$ \\
1 & 1 & 2 & 0 & 1 & $0 + 2(1) = 2$ \\
\bottomrule
\end{tabular}
\end{center}

In every case, $a_i + b_i = (a_i \oplus b_i) + 2(a_i \bitand b_i)$ holds identically.
Multiplying by $2^i$ and summing over all $i \ge 0$:
\begin{equation}
\sum_{i=0}^\infty (a_i + b_i) 2^i = \sum_{i=0}^\infty (a_i \oplus b_i) 2^i + 2 \sum_{i=0}^\infty (a_i \bitand b_i) 2^i.
\end{equation}
By linearity of summation, this yields $a + b = (a \oplus b) + 2(a \bitand b)$.
\end{proof}

\begin{corollary}[Carryless Addition Criterion]\label{cor:nim:carryless}
For any non-negative integers $a, b \in \mathbb{N}_0$,
\begin{equation}
a \oplus b = a + b \iff a \bitand b = 0.
\end{equation}
\end{corollary}

\begin{proof}
From Lemma~\ref{lem:nim:add_xor_identity}, $(a + b) - (a \oplus b) = 2(a \bitand b)$.
Since $a \bitand b \ge 0$, the difference vanishes if and only if $a \bitand b = 0$.
\end{proof}

Applying Corollary~\ref{cor:nim:carryless} to $a = n$ and $b = 2n$:
\begin{equation}\label{eq:nim:n_and_2n}
X(n, 2n, 3n) = 0 \iff n \bitand (2n) = 0.
\end{equation}


\section{The "No Consecutive Ones" Bitwise Characterization}

\subsection{Bit-Shift Analysis}
Let $n \in \mathbb{N}$ have binary representation $n = \sum_{i=0}^\infty n_i 2^i$ with $n_i \in \{0, 1\}$.
Multiplying $n$ by 2 corresponds to an arithmetic left-shift by 1 bit:
\begin{equation}
2n = \sum_{i=0}^\infty n_i 2^{i+1} = \sum_{j=1}^\infty n_{j-1} 2^j.
\end{equation}
Thus, the $i$-th bit of $2n$ is:
\begin{equation}
(2n)_0 = 0, \qquad (2n)_i = n_{i-1} \quad \text{for all } i \ge 1.
\end{equation}

\subsection{Bitwise AND Evaluation}
The $i$-th bit of the bitwise AND product $n \bitand (2n)$ is:
\begin{align}
(n \bitand 2n)_0 &= n_0 \cdot (2n)_0 = n_0 \cdot 0 = 0, \\
(n \bitand 2n)_i &= n_i \cdot (2n)_i = n_i \cdot n_{i-1} \quad \text{for all } i \ge 1.
\end{align}

Therefore, $n \bitand (2n) = 0$ if and only if every bit of the product is 0:
\begin{equation}
n \bitand (2n) = 0 \iff n_i \cdot n_{i-1} = 0 \quad \text{for all } i \ge 1.
\end{equation}

\begin{theorem}[Bit-Pattern Characterization]\label{thm:nim:no_consecutive_ones}
A positive integer $n$ satisfies $X(n, 2n, 3n) = 0$ if and only if the binary expansion of $n$ contains no adjacent ones (that is, $n$ contains no substring \texttt{11} in base 2).
\end{theorem}

\section{Combinatorial Enumeration via Fibonacci Numbers}

\subsection{Binary Strings Without Consecutive Ones}
Let $\mathcal{B}_k$ denote the set of binary words of length $k \ge 0$ over $\{0, 1\}$ containing no adjacent ones:
\begin{equation}
\mathcal{B}_k = \{ w \in \{0, 1\}^k : w \text{ does not contain the substring \texttt{11}} \}.
\end{equation}
Let $a_k = |\mathcal{B}_k|$.

\begin{itemize}
    \item \textbf{Base Cases:}
    \begin{itemize}
        \item $k = 0$: The empty word $\varepsilon$ contains no \texttt{11} $\implies a_0 = 1$.
        \item $k = 1$: The words \texttt{"0"} and \texttt{"1"} both contain no \texttt{11} $\implies a_1 = 2$.
    \end{itemize}
    
    \item \textbf{Recurrence for $k \ge 2$:}
    We partition $\mathcal{B}_k$ according to the prefix of each word $w \in \mathcal{B}_k$:
    \begin{enumerate}
        \item \textbf{Words starting with \texttt{0}:} Any such word is of the form \texttt{0}$u$, where $u \in \mathcal{B}_{k-1}$. There are $a_{k-1}$ such words.
        \item \textbf{Words starting with \texttt{1}:} Since \texttt{11} is forbidden, the second character must be \texttt{0}. Any such word is of the form \texttt{10}$v$, where $v \in \mathcal{B}_{k-2}$. There are $a_{k-2}$ such words.
    \end{enumerate}
    Since these two subsets are disjoint and partition $\mathcal{B}_k$:
    \begin{equation}\label{eq:nim:fib_recurrence}
    a_k = a_{k-1} + a_{k-2} \quad \text{for all } k \ge 2.
    \end{equation}
\end{itemize}

\subsection{Connection to the Fibonacci Sequence}
Recall the standard Fibonacci sequence $(F_m)_{m \ge 1}$:
\begin{equation}
F_1 = 1, \quad F_2 = 1, \quad F_m = F_{m-1} + F_{m-2} \quad (m \ge 3).
\end{equation}

\begin{lemma}\label{lem:nim:ak_fib}
For every integer $k \ge 0$,
\begin{equation}
a_k = F_{k+2}.
\end{equation}
\end{lemma}

\begin{proof}
We proceed by induction on $k$:
\begin{itemize}
    \item For $k = 0$: $a_0 = 1 = F_2 = F_{0+2}$.
    \item For $k = 1$: $a_1 = 2 = F_3 = F_{1+2}$.
    \item Inductive step: Assume $a_j = F_{j+2}$ holds for all $j < k$ ($k \ge 2$). Then:
    \begin{equation}
    a_k = a_{k-1} + a_{k-2} = F_{k+1} + F_k = F_{k+2}.
    \end{equation}
\end{itemize}
By mathematical induction, $a_k = F_{k+2}$ for all $k \ge 0$.
\end{proof}


\section{Counting Solutions on $[1, 2^K]$}

\subsection{Integers in the Half-Open Interval $[0, 2^K)$}
Every integer $n$ in the range $0 \le n < 2^K$ can be represented as a binary string of length exactly $K$ (by prepending leading zeros if necessary). Since leading zeros do not introduce \texttt{11} substrings, there is a natural bijection between $\{ n \in \mathbb{N}_0 : 0 \le n < 2^K \text{ and } n \bitand 2n = 0 \}$ and $\mathcal{B}_K$.
By Lemma~\ref{lem:nim:ak_fib}:
\begin{equation}
\#\{ n \in \mathbb{N}_0 : 0 \le n < 2^K \text{ and } n \bitand 2n = 0 \} = a_K = F_{K+2}.
\end{equation}

\subsection{Exclusion of $n = 0$}
The set above includes $n = 0$ ($0 \bitand 0 = 0$). Removing $n = 0$ leaves:
\begin{equation}
\#\{ n \in \mathbb{N} : 1 \le n < 2^K \text{ and } n \bitand 2n = 0 \} = F_{K+2} - 1.
\end{equation}

\subsection{Inclusion of the Upper Endpoint $n = 2^K$}
Consider the boundary integer $n = 2^K$:
\begin{itemize}
    \item The binary expansion of $2^K$ is $1\underbrace{00\dots00}_{K \text{ zeros}}$.
    \item This representation contains only a single 1-bit, and thus contains no adjacent ones.
    \item Directly checking: $2^K \bitand (2 \cdot 2^K) = 2^K \bitand 2^{K+1} = 0$.
\end{itemize}
Hence, $n = 2^K$ is always a valid solution.

\subsection{Total Solution Count $\mathcal{N}(2^K)$}
Summing the contributions:
\begin{equation}
\mathcal{N}(2^K) = (F_{K+2} - 1) + 1 = F_{K+2}.
\end{equation}

\begin{theorem}[Closed-Form Solution]\label{thm:nim:final_count}
For any integer $K \ge 1$, the number of positive integers $n \le 2^K$ such that $X(n, 2n, 3n) = 0$ is given by the $(K+2)$-th Fibonacci number:
\begin{equation}
\mathcal{N}(2^K) = F_{K+2}.
\end{equation}
\end{theorem}


\section{Exact Numerical Evaluation for $K = 30$}

Setting $K = 30$ in Theorem~\ref{thm:nim:final_count}, the total number of solutions is:
\begin{equation}
\mathcal{N}(2^{30}) = F_{30+2} = F_{32}.
\end{equation}

\subsection{Table of Fibonacci Numbers ($F_1$ to $F_{32}$)}

\begin{center}
\begin{tabular}{rr||rr||rr||rr}
\toprule
$m$ & $F_m$ & $m$ & $F_m$ & $m$ & $F_m$ & $m$ & $F_m$ \\
\midrule
1 & 1 & 9 & 34 & 17 & 1\,597 & 25 & 75\,025 \\
2 & 1 & 10 & 55 & 18 & 2\,584 & 26 & 121\,393 \\
3 & 2 & 11 & 89 & 19 & 4\,181 & 27 & 196\,418 \\
4 & 3 & 12 & 144 & 20 & 6\,765 & 28 & 317\,811 \\
5 & 5 & 13 & 233 & 21 & 10\,946 & 29 & 514\,229 \\
6 & 8 & 14 & 377 & 22 & 17\,711 & 30 & 832\,040 \\
7 & 13 & 15 & 610 & 23 & 28\,657 & 31 & 1\,346\,269 \\
8 & 21 & 16 & 987 & 24 & 46\,368 & \textbf{32} & \textbf{2\,178\,309} \\
\bottomrule
\end{tabular}
\end{center}

\subsection{Final Conclusion}
The number of positive integers $n \le 2^{30}$ for which $X(n, 2n, 3n) = 0$ is:
\begin{equation*}
\boxed{\mathbf{2\,178\,309}}
\end{equation*}

\end{document}
